% !PRINTED 784
\weberpage{784}
%<<CONT
$s$ in 1 übergehen lassen, so wird $\log Q_1(s)$ nach 4. unendlich, während die übrigen Glieder der linken Seite, $\log Q_2(s), \ldots \log Q_\mu(s)$, endlich bleiben. Da $R(s)$ gleichfalls endlich bleibt, so muss die Summe $\Sigma p^{-s}$ unendlich werden, und dies ist sicher nur dann möglich, wenn die Summe aus unendlich vielen Gliedern besteht. Hiermit ist bewiesen:

5. \gesperrt{In jeder der Zahlclassen $A$ (nach dem Modul $m$) sind unendlich viele Primzahlen enthalten.}

Man kann dies auch so ausdrücken, dass die Linearform $mx + a$, in der $m$ und $a$ ganze Zahlen ohne gemeinsamen Theiler sind, für unendlich viele ganzzahlige Werthe von $x$ eine Primzahl darstellt.

\begin{center}
\rule{4cm}{0.4pt}
\end{center}

% [The remainder of this leaf and the following leaf are blank in the source.]
